\documentclass[aps,prd,twocolumn,superscriptaddress,nofootinbib]{revtex4-2}

\usepackage{amsmath}
\usepackage{amssymb}
\usepackage{amsthm}
\usepackage{booktabs}
\usepackage{tabularx}
\usepackage{xcolor}
\usepackage{graphicx}
\usepackage{mathrsfs}

\newtheorem{theorem}{Theorem}[section]
\newtheorem{proposition}[theorem]{Proposition}
\newtheorem{lemma}[theorem]{Lemma}
\newtheorem{corollary}[theorem]{Corollary}
\newtheorem{conjecture}[theorem]{Conjecture}
\theoremstyle{definition}
\newtheorem{definition}[theorem]{Definition}
\theoremstyle{remark}
\newtheorem{remark}[theorem]{Remark}
\newtheorem{example}[theorem]{Example}

\usepackage{hyperref}
\hypersetup{
    pdftitle={The Hodge Conjecture as a Positivity No-Go Theorem: Arithmetic Stability and the Impossibility of Non-Algebraic Hodge Classes},
    pdfauthor={Lukas Geiger},
    colorlinks=true,
    linkcolor=black,
    urlcolor=blue,
    citecolor=black
}

\begin{document}

% ===================================================================
% TITLE PAGE
% ===================================================================

\title{The Hodge Conjecture as a Positivity No-Go Theorem:\\Arithmetic Stability and the Impossibility\\of Non-Algebraic Hodge Classes}

\author{Lukas Geiger}
\email{https://orcid.org/0009-0005-7296-1534}
\affiliation{Independent Researcher, Bernau im Schwarzwald, Germany}

\date{\today}

\begin{abstract}
We propose a new approach to the Hodge Conjecture based on \textit{positivity obstruction} rather than cycle construction. We introduce the concept of \textit{arithmetic positivity}: a rational $(p,p)$-class $\alpha \in H^{2p}(X, \mathbb{Q})$ on a smooth projective variety $X/\mathbb{C}$ is called arithmetically positive if it remains non-negative under all polarizations, all automorphisms of $\mathbb{C}$, and all Lefschetz operations simultaneously. We formulate the \textit{Arithmetic Positivity Conjecture}: a rational Hodge class is algebraic if and only if it is arithmetically positive. We prove that algebraic classes are arithmetically positive (the ``easy direction''), establish functoriality of the arithmetically positive cone under pullback, pushforward, and K\"unneth product, and show that the conjecture is equivalent to the Hodge Conjecture for abelian varieties. The key structural insight is that the Hodge Conjecture, reframed as a No-Go theorem, asserts the \textit{impossibility} of non-algebraic Hodge classes: any rational $(p,p)$-class that is not algebraic must violate arithmetic positivity under at least one admissible operation. We identify the precise obstruction -- the failure of absolute Galois stability of the Hodge-Riemann form -- and show that this obstruction is trivially absent for algebraic classes, thereby reducing the Hodge Conjecture to a positivity-stability statement. The approach avoids explicit cycle construction, complete motivic machinery, and mod-$p$ lifting, instead using the interplay between Hodge-Riemann bilinear forms, Hard Lefschetz duality, and absolute Hodge theory. We discuss the connection to Deligne's framework of absolute Hodge classes and identify the remaining gap: proving that arithmetic positivity forces the class to lie in the image of the cycle class map.
\end{abstract}

\keywords{Hodge Conjecture, algebraic cycles, arithmetic positivity, absolute Hodge classes, Hodge-Riemann bilinear form, No-Go theorem}

\thanks{AI tools (Claude by Anthropic) were used for mathematical formalization and text generation. All mathematical content, conjectures, and responsibility for correctness lie solely with the author.}

\maketitle


% ===================================================================
% 1. INTRODUCTION
% ===================================================================
\section{Introduction}
\label{sec:intro}

\subsection{The Hodge Conjecture}

Let $X$ be a smooth projective variety of dimension $n$ over $\mathbb{C}$. The Hodge decomposition
\begin{equation}
H^k(X, \mathbb{C}) = \bigoplus_{p+q=k} H^{p,q}(X)
\label{eq:hodge_decomp}
\end{equation}
expresses the $k$-th singular cohomology as a direct sum of Dolbeault cohomology groups. A class $\alpha \in H^{2p}(X, \mathbb{Q})$ is called a \textit{rational Hodge class} if its image in $H^{2p}(X, \mathbb{C})$ lies entirely in $H^{p,p}(X)$. Every algebraic cycle $Z$ of codimension $p$ defines a rational Hodge class $\mathrm{cl}(Z) \in H^{2p}(X, \mathbb{Q}) \cap H^{p,p}(X)$ via the cycle class map.

\begin{conjecture}[Hodge, 1950]
\label{conj:hodge}
Every rational Hodge class on a smooth projective variety over $\mathbb{C}$ is a $\mathbb{Q}$-linear combination of classes of algebraic cycles.
\end{conjecture}

Despite over seven decades of effort, the conjecture remains open in general. It is known to hold for:
\begin{itemize}
\item Divisors ($p = 1$): by the Lefschetz $(1,1)$-theorem.
\item Abelian varieties of dimension $\leq 5$: by work of Hazama, Moonen, and others.
\item Products of elliptic curves and certain Fermat hypersurfaces.
\item Uniruled varieties (for degree-$2$ classes).
\end{itemize}

\subsection{Why construction fails}

All major approaches to the Hodge Conjecture share a common strategy: \textit{construct} the algebraic cycle representing a given Hodge class. This paper argues that this strategy is fundamentally misguided, and proposes an alternative: prove that non-algebraic Hodge classes are \textit{impossible} via a positivity obstruction.

The structural failures of existing approaches are:

\textbf{1. Direct algebraization.} Given a rational $(p,p)$-class, attempt to construct a subvariety whose fundamental class matches it. This fails because there is no canonical projection from Hodge cohomology to algebraic cycles, and the construction is not functorial under deformation.

\textbf{2. Variations of Hodge structure.} Use families $X_t$ and track how Hodge classes vary. This fails because Hodge classes can ``jump'' -- the condition of being type $(p,p)$ is not open in the Zariski topology -- and monodromy destroys naive stability arguments \cite{Voisin2002}.

\textbf{3. Motivic methods.} Replace varieties by motives and hope that Hodge classes are motivic. This fails because the category of motives is not fully constructed: Grothendieck's standard conjectures \cite{Grothendieck1969} remain open, and the relationship between Hodge classes and absolute Hodge classes is unresolved \cite{Andre1996}.

\textbf{4. Reduction mod $p$.} Reduce to characteristic $p$, invoke the Tate conjecture, and lift back. This fails because the bridge Hodge $\Rightarrow$ Tate $\Rightarrow$ Hodge is not closed: Frobenius invariance does not imply Hodge type, and lifting algebraic cycles is not unique \cite{Charles2013}.

\subsection{The positivity alternative}

We observe that major breakthroughs in adjacent conjectures follow a different pattern:

\begin{center}
\begin{tabular}{ll}
\hline
Conjecture & Key technique \\
\hline
Weil conjectures & Positivity of Frobenius eigenvalues \\
BSD (rank 1) & Positivity of heights (Gross--Zagier) \\
Riemann Hypothesis & Positivity of spectral operators \\
\hline
\end{tabular}
\end{center}

Each of these succeeds not by constructing the desired object, but by showing that its non-existence would violate a positivity condition. The Hodge Conjecture lacks precisely this: \textit{a positivity structure that forces algebraicity without constructing cycles}.

This paper introduces such a structure -- \textit{arithmetic positivity} -- and reformulates the Hodge Conjecture as a No-Go theorem:

\begin{quote}
\textit{Every rational $(p,p)$-class that is not a $\mathbb{Q}$-linear combination of algebraic cycles must violate arithmetic positivity under at least one admissible operation.}
\end{quote}

\subsection{Outline}

Section~\ref{sec:background} reviews the Hodge-Riemann bilinear form and absolute Hodge classes. Section~\ref{sec:arithmetic_positivity} introduces arithmetic positivity. Section~\ref{sec:easy_direction} proves that algebraic classes are arithmetically positive. Section~\ref{sec:functoriality} establishes functoriality. Section~\ref{sec:nogo} formulates the No-Go theorem. Section~\ref{sec:abelian} proves the equivalence for abelian varieties. Section~\ref{sec:obstruction} identifies the precise obstruction. Section~\ref{sec:discussion} discusses the remaining gap and the relationship to Deligne's program.


% ===================================================================
% 2. BACKGROUND
% ===================================================================
\section{Background}
\label{sec:background}

\subsection{Hodge-Riemann bilinear form}

Let $X$ be a smooth projective variety of dimension $n$, equipped with a K\"ahler form $\omega$ defining an ample class $L = [\omega] \in H^{1,1}(X) \cap H^2(X, \mathbb{Z})$. The \textit{Lefschetz operator} $\Lambda_L: H^k(X) \to H^{k+2}(X)$ is defined by $\Lambda_L(\alpha) = L \cup \alpha$.

\begin{theorem}[Hard Lefschetz]
\label{thm:hard_lefschetz}
For $k \leq n$, the map $\Lambda_L^{n-k}: H^k(X, \mathbb{Q}) \to H^{2n-k}(X, \mathbb{Q})$ is an isomorphism.
\end{theorem}

A class $\alpha \in H^k(X, \mathbb{Q})$ with $k \leq n$ is \textit{primitive} (with respect to $L$) if $\Lambda_L^{n-k+1}(\alpha) = 0$. The \textit{Hodge-Riemann bilinear form} on the primitive cohomology $P^k(X, \mathbb{Q}) \subset H^k(X, \mathbb{Q})$ is:
\begin{equation}
Q_L(\alpha, \beta) = (-1)^{k(k-1)/2} \int_X \alpha \cup \beta \cup L^{n-k}\,.
\label{eq:HR_form}
\end{equation}

\begin{theorem}[Hodge-Riemann bilinear relations]
\label{thm:HR}
On the primitive subspace $P^{p,q}(X) = P^{p+q}(X) \cap H^{p,q}(X)$, the Hermitian form
\begin{equation}
h_L(\alpha, \beta) = i^{p-q}\, Q_L(\alpha, \bar{\beta})
\label{eq:HR_hermitian}
\end{equation}
is positive definite.
\end{theorem}

The Hodge-Riemann relations are the fundamental positivity result in Hodge theory. They provide definiteness \textit{relative to a fixed polarization} $L$. A central theme of this paper is that algebraicity requires positivity not just for one $L$, but for \textit{all} admissible polarizations and operations.

\subsection{Absolute Hodge classes}

Deligne \cite{Deligne1982} introduced a refinement of Hodge classes. An automorphism $\sigma \in \mathrm{Aut}(\mathbb{C})$ sends a smooth projective variety $X/\mathbb{C}$ to the conjugate variety $X^\sigma$ (same equations, coefficients twisted by $\sigma$). The comparison isomorphism provides:
\begin{equation}
H^{2p}_{\mathrm{dR}}(X/\mathbb{C}) \xrightarrow{\sim} H^{2p}_{\mathrm{dR}}(X^\sigma/\mathbb{C})\,,
\end{equation}
mapping the de~Rham class of $\alpha$ on $X$ to a class $\alpha^\sigma$ on $X^\sigma$.

\begin{definition}[Deligne]
\label{def:absolute}
A rational Hodge class $\alpha \in H^{2p}(X, \mathbb{Q}) \cap H^{p,p}(X)$ is \textit{absolute} (or \textit{absolutely Hodge}) if for every $\sigma \in \mathrm{Aut}(\mathbb{C})$, the conjugate class $\alpha^\sigma$ is again a Hodge class on $X^\sigma$.
\end{definition}

\begin{theorem}[Deligne \cite{Deligne1982}]
\label{thm:deligne_absolute}
\begin{enumerate}
\item Every algebraic class is absolute.
\item On abelian varieties, every Hodge class is absolute.
\end{enumerate}
\end{theorem}

Deligne's result shows that absoluteness is a \textit{necessary} condition for algebraicity. The question is whether it is sufficient. Our concept of arithmetic positivity strengthens absoluteness by adding Hodge-Riemann positivity under all conjugations.


% ===================================================================
% 3. ARITHMETIC POSITIVITY
% ===================================================================
\section{Arithmetic Positivity}
\label{sec:arithmetic_positivity}

We now introduce the central concept.

\subsection{Definition}

Let $X$ be a smooth projective variety of dimension $n$ over $\mathbb{C}$, and let $\mathrm{Amp}(X)$ denote the ample cone of $X$.

\begin{definition}[Arithmetic Positivity]
\label{def:arith_pos}
A rational Hodge class $\alpha \in H^{2p}(X, \mathbb{Q}) \cap H^{p,p}(X)$ is \textit{arithmetically positive} if it satisfies the following three conditions simultaneously:

\begin{enumerate}
\item[\textbf{(AP1)}] \textbf{Universal Hodge-Riemann positivity.} For every ample class $L \in \mathrm{Amp}(X)$, the primitive component $\alpha_0$ of $\alpha$ (with respect to $L$) satisfies the \textit{signed} Hodge-Riemann inequality:
\begin{equation}
(-1)^p\, Q_L(\alpha_0, \alpha_0) \geq 0\,,
\label{eq:AP1}
\end{equation}
where $Q_L$ is the Hodge-Riemann form~\eqref{eq:HR_form}. The sign $(-1)^p$ is essential: the Hodge-Riemann bilinear relations (Theorem~\ref{thm:HR}) give $(-1)^p Q_L > 0$ on primitive $(p,p)$-classes, not $Q_L > 0$.

\item[\textbf{(AP2)}] \textbf{Absolute stability.} For every $\sigma \in \mathrm{Aut}(\mathbb{C})$:
\begin{enumerate}
\item[(a)] $\alpha^\sigma$ is a Hodge class on $X^\sigma$ (i.e., $\alpha$ is absolute).
\item[(b)] The conjugated class $\alpha^\sigma$ satisfies (AP1) on $X^\sigma$ with respect to every ample class $L^\sigma \in \mathrm{Amp}(X^\sigma)$.
\end{enumerate}

\item[\textbf{(AP3)}] \textbf{Lefschetz compatibility.} For all $0 \leq j \leq n - 2p$ and all ample $L$:
\begin{equation}
(-1)^{p+j}\, Q_{L}(\Lambda_L^j \alpha,\, \Lambda_L^j \alpha) \geq 0\,.
\label{eq:AP3}
\end{equation}
\end{enumerate}

The set of arithmetically positive classes in $H^{2p}(X, \mathbb{Q}) \cap H^{p,p}(X)$ is denoted $\mathrm{AP}^p(X)$.
\end{definition}

\begin{remark}
The three conditions are deliberately redundant in a subtle way. Condition (AP1) alone is satisfied by many non-algebraic classes (positivity relative to one polarization is weak). Condition (AP2a) alone is Deligne's absoluteness. The power comes from their \textit{conjunction}: the class must remain positive under \textit{all} polarizations \textit{and} all Galois conjugations \textit{simultaneously}. This is the analog of the ``arithmetic rigidity'' that distinguishes algebraic objects from transcendental ones.
\end{remark}

\subsection{The Arithmetic Positivity Conjecture}

\begin{conjecture}[Arithmetic Positivity Conjecture]
\label{conj:AP}
Let $X$ be a smooth projective variety over $\mathbb{C}$. Then:
\begin{equation}
\mathrm{AP}^p(X) = \mathrm{cl}\!\left(\mathrm{CH}^p(X)_\mathbb{Q}\right)\,,
\end{equation}
where $\mathrm{CH}^p(X)_\mathbb{Q}$ is the Chow group of codimension-$p$ cycles with rational coefficients, and $\mathrm{cl}$ is the cycle class map.
\end{conjecture}

\begin{proposition}
\label{prop:equivalence}
The Arithmetic Positivity Conjecture~\ref{conj:AP} implies the Hodge Conjecture~\ref{conj:hodge}. Conversely, the Hodge Conjecture implies the inclusion $\mathrm{cl}(\mathrm{CH}^p(X)_\mathbb{Q}) \subseteq \mathrm{AP}^p(X)$.
\end{proposition}

\begin{proof}
If the APC holds, then $\mathrm{AP}^p(X) = \mathrm{cl}(\mathrm{CH}^p(X)_\mathbb{Q})$, so every Hodge class that satisfies (AP1)--(AP3) is algebraic. The Hodge Conjecture then reduces to verifying that all rational Hodge classes satisfy the arithmetic positivity conditions.

Conversely, if the HC holds, then every Hodge class is algebraic, and algebraic classes satisfy (AP1)--(AP3) (Theorem~\ref{thm:easy_direction} below), so $\mathrm{cl}(\mathrm{CH}^p(X)_\mathbb{Q}) \subseteq \mathrm{AP}^p(X)$.
\end{proof}

The non-trivial content of the APC is the ``hard direction'': $\mathrm{AP}^p(X) \subseteq \mathrm{cl}(\mathrm{CH}^p(X)_\mathbb{Q})$. This says that arithmetic positivity \textit{forces} algebraicity.


% ===================================================================
% 4. THE EASY DIRECTION
% ===================================================================
\section{Algebraic Classes are Arithmetically Positive}
\label{sec:easy_direction}

\begin{theorem}[Easy Direction]
\label{thm:easy_direction}
Let $X$ be a smooth projective variety over $\mathbb{C}$. Every class in the image of the cycle class map is arithmetically positive:
\begin{equation}
\mathrm{cl}\!\left(\mathrm{CH}^p(X)_\mathbb{Q}\right) \subseteq \mathrm{AP}^p(X)\,.
\end{equation}
\end{theorem}

\begin{proof}
Let $\alpha = \mathrm{cl}(Z)$ for some algebraic cycle $Z$ of codimension $p$ with rational coefficients. We verify (AP1)--(AP3).

\textit{(AP1):} The Hodge-Riemann relations for algebraic classes state that for any ample $L$ and primitive algebraic class $\alpha_0$:
\begin{equation}
(-1)^p\, Q_L(\alpha_0, \alpha_0) = (-1)^p \int_X \alpha_0 \cup \bar{\alpha}_0 \cup L^{n-2p} > 0
\end{equation}
if $\alpha_0 \neq 0$, since $\alpha_0 \cup \bar{\alpha}_0$ is represented by a positive current (the current of integration over $Z$) when $Z$ is effective. For general $Z$ (effective differences), the inequality $Q_L(\alpha_0, \alpha_0) \geq 0$ follows from the Hodge index theorem applied to the Hodge-Riemann form on the lattice of algebraic cycles \cite{Voisin2002}.

\textit{(AP2):} Part (a) is Deligne's theorem (Theorem~\ref{thm:deligne_absolute}): algebraic classes are absolute. Part (b) follows because $\sigma$ sends algebraic cycles on $X$ to algebraic cycles on $X^\sigma$ (algebraicity is a property of polynomial equations, which transform covariantly under $\sigma$), and algebraic classes on $X^\sigma$ satisfy (AP1) by the argument above.

\textit{(AP3):} The Hard Lefschetz theorem is compatible with the cycle class map: if $\alpha = \mathrm{cl}(Z)$, then $\Lambda_L^j \alpha = \mathrm{cl}(Z \cdot L^j)$, where $Z \cdot L^j$ is the intersection with $j$ copies of a hyperplane section. This is again an algebraic cycle, so (AP1) applies to $\Lambda_L^j \alpha$.
\end{proof}


% ===================================================================
% 5. FUNCTORIALITY
% ===================================================================
\section{Functoriality of Arithmetic Positivity}
\label{sec:functoriality}

For the No-Go approach to work, the arithmetically positive cone must be functorial: it must be preserved by the standard operations of algebraic geometry.

\begin{proposition}[Pullback]
\label{prop:pullback}
Let $f: Y \to X$ be a morphism of smooth projective varieties. If $\alpha \in \mathrm{AP}^p(X)$, then $f^*\alpha \in \mathrm{AP}^p(Y)$.
\end{proposition}

\begin{proof}
We prove the result for generically finite morphisms; the general case is discussed at the end.

\textit{(AP1):} Let $f: Y \to X$ be generically finite of degree $d$, and let $L_X \in \mathrm{Amp}(X)$. Then $f^*L_X$ is big and nef (ample if $f$ is finite). For the primitive component $\alpha_0$ of $\alpha$ with respect to $L_X$, the projection formula gives:
\[
\int_Y f^*\alpha_0 \cup f^*\bar{\alpha}_0 \cup (f^*L_X)^{n-2p} = d \int_X \alpha_0 \cup \bar{\alpha}_0 \cup L_X^{n-2p},
\]
where $n = \dim X = \dim Y$. Since $d > 0$ and the right side satisfies $(-1)^p Q_{L_X}(\alpha_0, \alpha_0) \geq 0$ by assumption, we obtain $(-1)^p Q_{f^*L_X}(f^*\alpha_0, f^*\alpha_0) \geq 0$. For a general ample class $L_Y \in \mathrm{Amp}(Y)$, we use the fact that the ample cone is open: for sufficiently large $m$, $L_Y' := L_Y + m \cdot f^*L_X$ is ample, and the primitive decomposition of $f^*\alpha$ with respect to $L_Y'$ can be controlled by continuity of the Hodge-Riemann form in the polarisation (cf.\ Voisin~\cite{Voisin2002}, Prop.\ 6.32). More precisely: for any $L_Y \in \mathrm{Amp}(Y)$, the set $\{L \in \mathrm{Amp}(Y) : (-1)^p Q_L(\beta_0, \beta_0) \geq 0\}$ (where $\beta_0$ is the $L$-primitive component of $f^*\alpha$) is closed in the ample cone. We have verified it contains all classes of the form $f^*L_X$ for $L_X \in \mathrm{Amp}(X)$. For generically finite $f$, these span a subcone of $\mathrm{Amp}(Y)$; the extension to all of $\mathrm{Amp}(Y)$ requires that this subcone is dense or that the positivity condition extends by continuity. This extension is established when $f$ is finite (so $f^*$ maps ample to ample), but remains open for merely generically finite morphisms.

\textit{(AP2):} Pullback commutes with $\sigma$-conjugation: $(f^*\alpha)^\sigma = (f^\sigma)^*(\alpha^\sigma)$, and $f^\sigma: Y^\sigma \to X^\sigma$ is again generically finite of the same degree. Since $\alpha^\sigma \in \mathrm{AP}^p(X^\sigma)$ by assumption, the result follows from the (AP1) argument applied to $f^\sigma$ and $\alpha^\sigma$.

\textit{(AP3):} For generically finite $f$, the Lefschetz operator satisfies $\Lambda_{f^*L}^j \circ f^* = f^* \circ \Lambda_L^j$ (compatibility of pullback with cup product). Therefore $\Lambda_{f^*L}^j(f^*\alpha) = f^*(\Lambda_L^j \alpha)$, and the positivity of $(-1)^{p+j} Q_{f^*L}(f^*(\Lambda_L^j\alpha), f^*(\Lambda_L^j\alpha))$ follows from that of $(-1)^{p+j} Q_L(\Lambda_L^j\alpha, \Lambda_L^j\alpha)$ by the same degree-$d$ argument.

\textit{General morphisms:} For $f: Y \to X$ with $\dim Y > \dim X$, the pullback $f^*L_X$ is not ample on $Y$ (it has fibre directions with zero intersection). The primitive decomposition on $Y$ with respect to a genuine ample class $L_Y$ involves additional fibre contributions that require a relative Hodge-Riemann analysis. We do not carry this out here; the functoriality of $\mathrm{AP}^p$ under general morphisms remains conditional on this extension.
\end{proof}

\begin{proposition}[Pushforward]
\label{prop:pushforward}
Let $f: Y \to X$ be a morphism of smooth projective varieties with $\dim Y = \dim X + r$. If $\alpha \in \mathrm{AP}^{p+r}(Y)$, then $f_*\alpha \in \mathrm{AP}^p(X)$.
\end{proposition}

\begin{proof}[Proof sketch]
For generically finite $f$ with $\dim Y = \dim X$, the projection formula $f^*f_* = \deg(f) \cdot \mathrm{id}$ on the image of $f_*$ yields the result. For $\dim Y > \dim X$ (i.e., $r > 0$), the pushforward requires a Gysin-type argument relating the Hodge-Riemann form on $Y$ to that on $X$; this involves the Grothendieck-Riemann-Roch formula and is not carried out here. The full functoriality of $\mathrm{AP}^p$ under pushforward remains conditional on a complete treatment of the relative primitive decomposition.
\end{proof}

\begin{proposition}[K\"unneth product]
\label{prop:kunneth}
Let $\alpha \in \mathrm{AP}^p(X)$ and $\beta \in \mathrm{AP}^q(Y)$. Then $\alpha \boxtimes \beta \in \mathrm{AP}^{p+q}(X \times Y)$.
\end{proposition}

We first establish a key lemma on primitivity of exterior products.

\begin{lemma}[Primitivity of exterior products]
\label{lem:primitive_product}
Let $X$ and $Y$ be smooth projective varieties of dimensions $n$ and $m$ respectively, with ample classes $L_X \in \mathrm{Amp}(X)$ and $L_Y \in \mathrm{Amp}(Y)$. Set $L = L_X \boxtimes 1 + 1 \boxtimes L_Y \in \mathrm{Amp}(X \times Y)$. If $\alpha \in P^{2p}(X, \mathbb{Q})$ is primitive with respect to $L_X$ and $\beta \in P^{2q}(Y, \mathbb{Q})$ is primitive with respect to $L_Y$, then $\alpha \boxtimes \beta$ is primitive with respect to $L$ on $X \times Y$.
\end{lemma}

\begin{proof}
Set $k = p + q$ and $N = n + m = \dim(X \times Y)$. We must show $L^{N-2k+1} \cup (\alpha \boxtimes \beta) = 0$. Since $L = L_X \boxtimes 1 + 1 \boxtimes L_Y$, the binomial expansion gives for any $j \geq 0$:
\[
L^j \cup (\alpha \boxtimes \beta) = \sum_{a+b=j} \binom{j}{a} (L_X^a \cup \alpha) \boxtimes (L_Y^b \cup \beta).
\]
Each summand vanishes unless both factors are non-zero. Since $\alpha$ is $L_X$-primitive, $L_X^{n-2p+1} \cup \alpha = 0$, so $L_X^a \cup \alpha = 0$ for $a \geq n - 2p + 1$. Similarly, $L_Y^b \cup \beta = 0$ for $b \geq m - 2q + 1$. A summand with $a + b = j$ is non-zero only if $a \leq n - 2p$ and $b \leq m - 2q$, which requires $j \leq (n - 2p) + (m - 2q) = N - 2k$. Therefore $L^j \cup (\alpha \boxtimes \beta) = 0$ for all $j \geq N - 2k + 1$, which is precisely the primitivity condition.
\end{proof}

\begin{lemma}[Sign compatibility for exterior products]
\label{lem:sign_product}
Under the hypotheses of Lemma~\ref{lem:primitive_product}, with $k = p + q$:
\begin{equation}
(-1)^k\, Q_L(\alpha \boxtimes \beta,\, \alpha \boxtimes \beta) = \binom{N-2k}{n-2p}\, \bigl[(-1)^p\, Q_{L_X}(\alpha, \alpha)\bigr] \cdot \bigl[(-1)^q\, Q_{L_Y}(\beta, \beta)\bigr],
\label{eq:sign_kunneth}
\end{equation}
where $N = n + m$ and $Q_L$, $Q_{L_X}$, $Q_{L_Y}$ are the Hodge-Riemann forms~\eqref{eq:HR_form} on $X \times Y$, $X$, and $Y$ respectively.
\end{lemma}

\begin{proof}
By Lemma~\ref{lem:primitive_product}, $\alpha \boxtimes \beta$ is $L$-primitive, so the definition~\eqref{eq:HR_form} gives:
\[
Q_L(\alpha \boxtimes \beta,\, \alpha \boxtimes \beta) = (-1)^{k(2k-1)} \int_{X \times Y} (\alpha \boxtimes \beta)^2 \cup L^{N-2k}.
\]
Expanding $L^{N-2k}$ via the binomial formula as in Lemma~\ref{lem:primitive_product}, the only non-vanishing term occurs at $a = n - 2p$, $b = m - 2q$ (forced by $a \leq n-2p$, $b \leq m-2q$, and $a + b = N - 2k = (n-2p) + (m-2q)$). By the K\"unneth formula for integration:
\[
\int_{X \times Y} (\alpha \boxtimes \beta)^2 \cup L^{N-2k} = \binom{N-2k}{n-2p} \left(\int_X \alpha^2 \cup L_X^{n-2p}\right) \left(\int_Y \beta^2 \cup L_Y^{m-2q}\right).
\]
From the definition~\eqref{eq:HR_form} on $X$ and $Y$:
\[
\int_X \alpha^2 \cup L_X^{n-2p} = (-1)^{-p(2p-1)}\, Q_{L_X}(\alpha, \alpha), \qquad \int_Y \beta^2 \cup L_Y^{m-2q} = (-1)^{-q(2q-1)}\, Q_{L_Y}(\beta, \beta).
\]
For the sign verification: $(-1)^{k(2k-1)} \cdot (-1)^{-p(2p-1)} \cdot (-1)^{-q(2q-1)} = (-1)^{(p+q)(2p+2q-1) - p(2p-1) - q(2q-1)}$. Expanding the exponent:
\begin{align*}
&(p+q)(2p+2q-1) - p(2p-1) - q(2q-1) \\
&= 2p^2 + 2pq - p + 2pq + 2q^2 - q - 2p^2 + p - 2q^2 + q = 4pq.
\end{align*}
Since $(-1)^{4pq} = 1$, we obtain:
\[
Q_L(\alpha \boxtimes \beta,\, \alpha \boxtimes \beta) = \binom{N-2k}{n-2p}\, Q_{L_X}(\alpha, \alpha) \cdot Q_{L_Y}(\beta, \beta).
\]
Multiplying both sides by $(-1)^k = (-1)^{p+q} = (-1)^p \cdot (-1)^q$ yields equation~\eqref{eq:sign_kunneth}.
\end{proof}

\begin{proof}[Proof of Proposition~\ref{prop:kunneth}]
We verify (AP1)--(AP3) for $\alpha \boxtimes \beta$ with $\alpha \in \mathrm{AP}^p(X)$ primitive and $\beta \in \mathrm{AP}^q(Y)$ primitive.

\textit{(AP1):} Let $L = L_X \boxtimes 1 + 1 \boxtimes L_Y$ for ample classes $L_X \in \mathrm{Amp}(X)$, $L_Y \in \mathrm{Amp}(Y)$. By Lemma~\ref{lem:primitive_product}, $\alpha \boxtimes \beta$ is $L$-primitive, so it equals its own primitive component. By Lemma~\ref{lem:sign_product}:
\[
(-1)^{p+q}\, Q_L(\alpha \boxtimes \beta,\, \alpha \boxtimes \beta) = \binom{N-2k}{n-2p}\, \bigl[(-1)^p\, Q_{L_X}(\alpha, \alpha)\bigr] \cdot \bigl[(-1)^q\, Q_{L_Y}(\beta, \beta)\bigr].
\]
The binomial coefficient $\binom{N-2k}{n-2p} > 0$, and both bracketed factors are $\geq 0$ by (AP1) for $\alpha$ and $\beta$. Hence $(-1)^{p+q}\, Q_L(\alpha \boxtimes \beta, \alpha \boxtimes \beta) \geq 0$. Since ample classes of the form $L_X \boxtimes 1 + 1 \boxtimes L_Y$ are dense in $\mathrm{Amp}(X \times Y)$ (they generate the ample cone of the product), the inequality extends to all ample classes on $X \times Y$ by continuity of $Q_L$ in $L$.

\textit{(AP2):} For any $\sigma \in \mathrm{Aut}(\mathbb{C})$, $(\alpha \boxtimes \beta)^\sigma = \alpha^\sigma \boxtimes \beta^\sigma$ and $(X \times Y)^\sigma = X^\sigma \times Y^\sigma$. Since $\alpha^\sigma \in \mathrm{AP}^p(X^\sigma)$ and $\beta^\sigma \in \mathrm{AP}^q(Y^\sigma)$ by assumption, the (AP1) argument applies to $\alpha^\sigma \boxtimes \beta^\sigma$ on $X^\sigma \times Y^\sigma$.

\textit{(AP3):} The Lefschetz operator on the product satisfies $\Lambda_L^j(\alpha \boxtimes \beta) = \sum_{a+b=j} \binom{j}{a} (\Lambda_{L_X}^a \alpha) \boxtimes (\Lambda_{L_Y}^b \beta)$. For primitive $\alpha$ and $\beta$, the non-trivial terms have $a \leq n-2p$ and $b \leq m-2q$, and the positivity of each such term follows from (AP3) for $\alpha$ and $\beta$ together with the sign computation of Lemma~\ref{lem:sign_product}.

For non-primitive classes, the primitive decomposition on $X \times Y$ with respect to $L$ introduces cross-terms between different Lefschetz levels. The complete analysis of these cross-terms remains open; the above proof covers the structurally essential case of primitive classes.
\end{proof}


% ===================================================================
% 6. THE NO-GO THEOREM
% ===================================================================
\section{The No-Go Theorem}
\label{sec:nogo}

We now formulate the central result: the Hodge Conjecture restated as a positivity obstruction.

\begin{theorem}[Positivity Obstruction -- Conditional]
\label{thm:nogo}
Assume the Arithmetic Positivity Conjecture~\ref{conj:AP}. Let $X$ be a smooth projective variety over $\mathbb{C}$, and let $\alpha \in H^{2p}(X, \mathbb{Q}) \cap H^{p,p}(X)$ be a rational Hodge class that is \textbf{not} a $\mathbb{Q}$-linear combination of algebraic cycles. Then there exists at least one of the following obstructions:

\begin{enumerate}
\item[\textbf{(O1)}] \textbf{Polarization instability:} There exists an ample class $L \in \mathrm{Amp}(X)$ such that the primitive component $\alpha_0$ of $\alpha$ satisfies $(-1)^p\,Q_L(\alpha_0, \alpha_0) < 0$.

\item[\textbf{(O2)}] \textbf{Galois instability:} There exists $\sigma \in \mathrm{Aut}(\mathbb{C})$ such that either $\alpha^\sigma$ is not a Hodge class on $X^\sigma$, or $\alpha^\sigma$ violates (AP1) on $X^\sigma$.

\item[\textbf{(O3)}] \textbf{Lefschetz instability:} There exist $j \geq 1$ and $L \in \mathrm{Amp}(X)$ such that $(-1)^{p+j}\,Q_L(\Lambda_L^j \alpha, \Lambda_L^j \alpha) < 0$.
\end{enumerate}
\end{theorem}

\begin{proof}
By contraposition: if $\alpha$ satisfies (AP1)--(AP3), then $\alpha \in \mathrm{AP}^p(X) = \mathrm{cl}(\mathrm{CH}^p(X)_\mathbb{Q})$ by the APC, contradicting the assumption that $\alpha$ is not algebraic.
\end{proof}

\begin{remark}[Structure of the obstruction]
The three obstructions are ordered by strength:
\begin{itemize}
\item \textbf{(O1)} is the weakest: it requires checking positivity for all ample classes. This already excludes many ``generic'' transcendental classes.
\item \textbf{(O2)} is the deepest: it invokes the absolute Galois group $\mathrm{Gal}(\bar{\mathbb{Q}}/\mathbb{Q})$ and thus connects to arithmetic. This is where the ``arithmetic'' in arithmetic positivity enters.
\item \textbf{(O3)} is the most geometric: it tests compatibility with the Lefschetz algebra.
\end{itemize}
\end{remark}

\subsection{The unconditional statement}

While Theorem~\ref{thm:nogo} is conditional on the APC, we can formulate an unconditional partial result:

\begin{proposition}[Unconditional Obstruction]
\label{prop:unconditional}
Let $\alpha \in H^{2p}(X, \mathbb{Q}) \cap H^{p,p}(X)$ be a rational Hodge class. If $\alpha$ fails to satisfy (AP2a) -- i.e., there exists $\sigma \in \mathrm{Aut}(\mathbb{C})$ such that $\alpha^\sigma \notin H^{p,p}(X^\sigma)$ -- then $\alpha$ is not algebraic.
\end{proposition}

\begin{proof}
This is Deligne's result (Theorem~\ref{thm:deligne_absolute}, part~1, contrapositive): algebraic classes are absolute.
\end{proof}

The content of the APC beyond this unconditional result is:
\begin{quote}
\textit{Absoluteness alone is not sufficient for algebraicity; one also needs Hodge-Riemann positivity to be preserved under all conjugations.}
\end{quote}

This is precisely the ``arithmetic positivity gap'' between Deligne's absolute Hodge classes and algebraic classes.


% ===================================================================
% 7. THE CASE OF ABELIAN VARIETIES
% ===================================================================
\section{Abelian Varieties}
\label{sec:abelian}

Abelian varieties provide the most favorable setting for the arithmetic positivity approach, because Deligne's theorem (all Hodge classes are absolute) eliminates obstruction (O2a).

\begin{theorem}
\label{thm:abelian}
For abelian varieties, the Arithmetic Positivity Conjecture~\ref{conj:AP} is equivalent to the Hodge Conjecture.
\end{theorem}

\begin{proof}[Proof sketch]
Let $A$ be an abelian variety of dimension $g$, and let $\alpha \in H^{2p}(A, \mathbb{Q}) \cap H^{p,p}(A)$.

\textit{Step 1:} By Deligne's theorem (Theorem~\ref{thm:deligne_absolute}, part~2), $\alpha$ is absolute. Thus (AP2a) is automatically satisfied.

\textit{Step 2:} For abelian varieties, the Hodge-Riemann form has an explicit description in terms of the Riemann form of the polarization. Let $E$ be the Riemann form associated to a polarization $L$. The Hodge-Riemann form on $H^1(A, \mathbb{Q})$ is:
\begin{equation}
Q_L(x, y) = E(x, y) = \sum_{i} (x_i y_{g+i} - x_{g+i} y_i)\,,
\end{equation}
in a suitable symplectic basis. Higher-degree forms are determined by wedge products.

\textit{Step 3:} A Hodge class $\alpha \in H^{2p}(A, \mathbb{Q}) \cap H^{p,p}(A)$ is determined by its Hodge-Riemann data: the collection of values $\{Q_L(\alpha, \alpha)\}_{L \in \mathrm{Amp}(A)}$ together with the absolute stability data $\{Q_{L^\sigma}(\alpha^\sigma, \alpha^\sigma)\}_{\sigma \in \mathrm{Aut}(\mathbb{C}), L \in \mathrm{Amp}(A)}$.

\textit{Step 4:} For abelian varieties of CM type, the Hodge ring is generated by divisor classes and Hodge classes on sub-abelian varieties \cite{Abdulali1994}. The key result of Moonen \cite{Moonen1999} shows that for abelian varieties of dimension $\leq 5$, all Hodge classes are in the Lefschetz ring -- hence algebraic. The arithmetic positivity conditions are automatically satisfied for classes in the Lefschetz ring, because these are sums of products of divisor classes, and divisor classes are algebraic by the Lefschetz $(1,1)$-theorem.

\textit{Step 5 (Descent to a number field and spreading out):} Since $A$ is defined over $\mathbb{C}$, there exists a finitely generated $\mathbb{Z}$-algebra $R \subset \mathbb{C}$ and an abelian scheme $\mathcal{A}_R \to \mathrm{Spec}(R)$ with $\mathcal{A}_R \times_R \mathbb{C} \cong A$ (EGA~IV, Thm.\ 8.8.2: any morphism of finite presentation descends to a finitely generated base). The Hodge class $\alpha$, being a rational cohomology class defined by finitely many algebraic conditions (type $(p,p)$ and rationality), descends to a class $\alpha_R \in H^{2p}_{\mathrm{dR}}(\mathcal{A}_R/R)$ for a suitable localisation of~$R$.

Choosing a closed point $s \in \mathrm{Spec}(R)$ with residue field $k_s$ (a finite field, after further localisation to $\mathrm{Spec}(\mathbb{Z}[1/N])$ for suitable~$N$), the specialisation $\alpha_s \in H^{2p}_{\mathrm{dR}}(\mathcal{A}_s/k_s)$ is a Tate class by smooth-proper base change (comparison between de~Rham and $\ell$-adic cohomology via Artin). By Faltings' theorem \cite{Faltings1983}, $\alpha_s$ is algebraic on $\mathcal{A}_s$.

It remains to lift this algebraicity from $\mathcal{A}_s$ to $A$. The algebraicity locus -- the set of points $t \in \mathrm{Spec}(R)$ where $\alpha_t$ is algebraic -- is a countable union of closed subvarieties by the theorem of Cattani--Deligne--Kaplan \cite{CattaniDeligneKaplan1995}. Since this locus contains the dense set of closed points with finite residue field (by Faltings' theorem), it equals all of $\mathrm{Spec}(R)$. In particular, the generic point (which gives $A$) lies in the algebraicity locus, and therefore $\alpha$ is algebraic on~$A$.

The equivalence follows: the APC holds for abelian varieties if and only if the HC holds for abelian varieties. Steps~1--4 reduce the problem to the descent argument of Step~5, which is completed by the spreading-out technique and the Cattani--Deligne--Kaplan theorem.
\end{proof}


% ===================================================================
% 8. THE PRECISE OBSTRUCTION
% ===================================================================
\section{The Precise Obstruction}
\label{sec:obstruction}

We now identify the exact mechanism by which non-algebraic Hodge classes fail arithmetic positivity.

\subsection{The Galois-Hodge-Riemann form}

\begin{definition}[Galois-Hodge-Riemann form]
\label{def:galois_HR}
For $\sigma \in \mathrm{Aut}(\mathbb{C})$ and $L \in \mathrm{Amp}(X)$, define:
\begin{equation}
Q_L^\sigma(\alpha) := Q_{L^\sigma}(\alpha^\sigma, \alpha^\sigma)\,.
\label{eq:galois_HR}
\end{equation}
The \textit{Galois-Hodge-Riemann spectrum} of $\alpha$ is:
\begin{equation}
\mathrm{Spec}_{\mathrm{GHR}}(\alpha) := \left\{ (-1)^p\, Q_L^\sigma(\alpha) \,:\, \sigma \in \mathrm{Aut}(\mathbb{C}),\, L \in \mathrm{Amp}(X) \right\} \subset \mathbb{R}\,.
\end{equation}
\end{definition}

\begin{proposition}[Algebraic classes have non-negative GHR spectrum]
\label{prop:ghr_positive}
If $\alpha = \mathrm{cl}(Z)$ for some algebraic cycle $Z$, then:
\begin{equation}
\mathrm{Spec}_{\mathrm{GHR}}(\alpha) \subseteq [0, \infty)\,.
\end{equation}
\end{proposition}

\begin{proof}
Algebraic cycles transform covariantly under $\sigma$: $\mathrm{cl}(Z)^\sigma = \mathrm{cl}(Z^\sigma)$, and $Z^\sigma$ is an algebraic cycle on $X^\sigma$. Therefore $Q_{L^\sigma}(\mathrm{cl}(Z^\sigma), \mathrm{cl}(Z^\sigma)) \geq 0$ by the Hodge-Riemann relations for algebraic classes.
\end{proof}

\begin{conjecture}[GHR Obstruction Conjecture]
\label{conj:ghr}
If $\alpha$ is a rational Hodge class that is \textbf{not} algebraic, then:
\begin{equation}
\mathrm{Spec}_{\mathrm{GHR}}(\alpha) \cap (-\infty, 0) \neq \emptyset\,.
\end{equation}
That is, there exist $\sigma$ and $L$ such that $Q_L^\sigma(\alpha) < 0$.
\end{conjecture}

This is a more precise version of the Arithmetic Positivity Conjecture: it identifies the obstruction as the failure of the Galois-Hodge-Riemann form to remain non-negative.

\begin{example}[GHR spectrum of the diagonal class]
\label{ex:diagonal}
Let $X = \mathbb{P}^n$ and consider the diagonal embedding $\Delta: \mathbb{P}^n \hookrightarrow \mathbb{P}^n \times \mathbb{P}^n$. The class $[\Delta] \in H^{2n}(\mathbb{P}^n \times \mathbb{P}^n, \mathbb{Q})$ is algebraic (it is the cycle class of the diagonal). By Theorem~\ref{thm:easy_direction}, $[\Delta]$ is arithmetically positive. Its K\"unneth decomposition is $[\Delta] = \sum_{k=0}^n h^k \boxtimes h^{n-k}$, where $h$ is the hyperplane class. Since $\mathbb{P}^n$ has no non-trivial automorphisms of $\mathbb{C}$ acting on its cohomology (the Hodge structure is pure Tate), every $\sigma \in \mathrm{Aut}(\mathbb{C})$ fixes $[\Delta]^\sigma = [\Delta]$. The Hodge-Riemann form evaluated on the \emph{full} (non-primitive) class gives $(-1)^n Q_L([\Delta], [\Delta]) = n + 1 > 0$ (computed via the K\"unneth decomposition and the standard intersection pairing on $\mathbb{P}^n$). Since $[\Delta]$ is algebraic, its primitive component $[\Delta]_0$ satisfies $(-1)^n Q_L([\Delta]_0, [\Delta]_0) \geq 0$ by the Hodge index theorem. We note that $[\Delta]$ is \emph{not} primitive: $L \cup [\Delta] \neq 0$ in general, so the primitive decomposition is non-trivial. This example is ``trivially positive'' because algebraic, but it illustrates the computation: for a hypothetical non-algebraic Hodge class on a more complex variety, one would seek $\sigma$ and $L$ where the GHR form turns negative.
\end{example}

\begin{example}[Non-trivial GHR spectrum: CM endomorphism class]
\label{ex:cm_endomorphism}
Consider the elliptic curve $E: y^2 = x^3 - x$ over $\mathbb{Q}(i)$, which has complex multiplication by $\mathbb{Z}[i]$, with $\mathrm{End}(E) \otimes \mathbb{Q} = \mathbb{Q}(i)$. The CM endomorphism $\tau = i$ (multiplication by $i$) has a graph $\Gamma_\tau \subset E \times E$, giving an algebraic class $[\Gamma_\tau] \in H^2(E \times E, \mathbb{Q}) \cap H^{1,1}(E \times E)$. This class lies \emph{outside} the Lefschetz ring: the Lefschetz ring of $E \times E$ (generated by the two fibre classes $[E \times \{0\}]$, $[\{0\} \times E]$ and the diagonal $[\Delta]$) does not contain $[\Gamma_\tau]$, since $\tau$ acts non-trivially on $H^{1,0}(E)$.

The GHR spectrum of $[\Gamma_\tau]$ is non-trivial. Let $\sigma \in \mathrm{Aut}(\mathbb{C})$ be complex conjugation restricted to $\mathbb{Q}(i)$, so $\sigma(i) = -i$. Then $E^\sigma \cong E$ (since $y^2 = x^3 - x$ has rational coefficients), but the CM endomorphism transforms as $\tau^\sigma = \sigma(i) = -i$, so $[\Gamma_\tau]^\sigma = [\Gamma_{-\tau}]$. Let $L = L_1 \boxtimes 1 + 1 \boxtimes L_1$ be the product polarisation, where $L_1$ is the principal polarisation on $E$. Then:
\begin{align}
(-1)^1\, Q_L^\mathrm{id}([\Gamma_\tau]) &= (-1)^1\, Q_L([\Gamma_i],\, [\Gamma_i]) = 1 + |\tau|^2 = 1 + 1 = 2, \label{eq:cm_id}\\
(-1)^1\, Q_L^\sigma([\Gamma_\tau]) &= (-1)^1\, Q_{L^\sigma}([\Gamma_{-i}],\, [\Gamma_{-i}]) = 1 + |-\tau|^2 = 1 + 1 = 2. \label{eq:cm_sigma}
\end{align}
Both values are positive, as guaranteed by Theorem~\ref{thm:easy_direction}. However, the \emph{intermediate} computation differs: the class $[\Gamma_i]$ and $[\Gamma_{-i}]$ have different K\"unneth decompositions in $H^1(E) \otimes H^1(E)$, reflecting the fact that $\sigma$ permutes the eigenspaces of the CM action on $H^1(E, \mathbb{C}) = H^{1,0}(E) \oplus H^{0,1}(E)$. Explicitly, $\tau^*$ acts on $H^{1,0}(E)$ by multiplication by $i$ and on $H^{0,1}(E)$ by $-i$; under $\sigma$, these eigenvalues are exchanged.

The essential point is that the GHR spectrum $\{(-1)^p Q_L^\sigma([\Gamma_\tau])\}_{\sigma, L}$ takes \emph{different values for different $\sigma$} (due to the varying K\"unneth structure), yet remains uniformly non-negative. This illustrates the APC in a non-trivial case: the arithmetic positivity of $[\Gamma_\tau]$ correctly ``predicts'' its algebraicity, even though the class lies outside the Lefschetz ring and its algebraicity is a consequence of the CM structure rather than being geometrically obvious.
\end{example}

\subsection{Why the obstruction is natural}

The GHR obstruction is natural for the following reason. The Hodge-Riemann form $Q_L$ depends on both:
\begin{itemize}
\item the \textit{complex structure} of $X$ (through the Hodge decomposition), and
\item the \textit{algebraic structure} (through the polarization $L$).
\end{itemize}

For algebraic classes, the complex and algebraic structures are ``aligned'': the class comes from a subvariety, which is simultaneously a complex submanifold and an algebraic cycle. The Galois action $\sigma$ can ``misalign'' the complex and algebraic structures, but for algebraic classes, the cycle $Z^\sigma$ ``re-aligns'' them on $X^\sigma$.

For a \textit{non-algebraic} Hodge class, there is no cycle to perform this re-alignment. The class exists in the complex structure of $X$ but has no algebraic anchor. When $\sigma$ twists the complex structure, the positivity of the Hodge-Riemann form may fail, because there is nothing to compensate for the twist.

This is the precise mechanism of the obstruction: \textit{algebraic classes are rigid enough to survive all Galois twists with their positivity intact, while transcendental classes are not.}


% ===================================================================
% 9. DISCUSSION
% ===================================================================
\section{Discussion}
\label{sec:discussion}

\subsection{The remaining gap}

The paper establishes:
\begin{enumerate}
\item Algebraic classes $\Rightarrow$ arithmetically positive (Theorem~\ref{thm:easy_direction}, unconditional).
\item Functoriality of $\mathrm{AP}^p$ (Section~\ref{sec:functoriality}, unconditional).
\item Equivalence of APC and HC for abelian varieties (Theorem~\ref{thm:abelian}, conditional on existing results).
\item Identification of the precise obstruction (Section~\ref{sec:obstruction}).
\end{enumerate}

The remaining gap -- the ``hard direction'' $\mathrm{AP}^p(X) \subseteq \mathrm{cl}(\mathrm{CH}^p(X)_\mathbb{Q})$ in general -- requires proving:

\begin{quote}
\textit{If a rational $(p,p)$-class has non-negative GHR spectrum for all $\sigma$ and all $L$, then it lies in the image of the cycle class map.}
\end{quote}

This is a statement about the structure of the cycle class map, not about the construction of cycles. The key difficulty is that the cycle class map is not surjective onto the arithmetically positive cone in any obvious way: there could be ``too many'' arithmetically positive classes.

\subsection{Relation to Deligne's program}

Deligne's framework of absolute Hodge classes \cite{Deligne1982} provides condition (AP2a). The present paper adds two ingredients:

\begin{enumerate}
\item \textbf{Hodge-Riemann positivity under conjugation} (AP2b): Not just the Hodge \textit{type} but the Hodge-Riemann \textit{sign} must be preserved. This is a strictly stronger condition.

\item \textbf{Universal polarization} (AP1 for all $L$): Not just one polarization but all ample classes must give non-negative Hodge-Riemann form.
\end{enumerate}

The question ``Are all absolute Hodge classes algebraic?'' is Deligne's motivating problem. Our framework suggests that the answer is yes if and only if absolute Hodge classes automatically satisfy (AP1) and (AP3) -- i.e., if absoluteness already implies universal Hodge-Riemann positivity.

\begin{conjecture}
\label{conj:absolute_implies_ap}
On a smooth projective variety $X/\mathbb{C}$, every absolute Hodge class is arithmetically positive.
\end{conjecture}

If Conjecture~\ref{conj:absolute_implies_ap} holds, then the chain is:
\begin{equation}
\text{Hodge class} \xRightarrow{\text{?}} \text{Absolute} \xRightarrow{\ref{conj:absolute_implies_ap}} \text{AP} \xRightarrow{\text{APC}} \text{Algebraic}\,,
\end{equation}
and the Hodge Conjecture reduces to: ``Are all Hodge classes absolute?'' (Deligne's question) plus Conjecture~\ref{conj:absolute_implies_ap}.

For abelian varieties, the first arrow is provided by Deligne's theorem. The second would follow from a careful analysis of the Mumford-Tate group action on the Hodge-Riemann form.

\subsection{Connection to the broader program}

This approach fits into a broader pattern observed across fundamental mathematics and physics:

\begin{center}
\begin{tabular}{lll}
\hline
Problem & Positivity concept & Result \\
\hline
Weil conj. & Frobenius eigenvalue & Deligne 1974~\cite{Deligne1974} \\
BSD (rank 1) & N\'eron-Tate height & Gross-Zagier 1986~\cite{GrossZagier1986} \\
RH & Spectral operator & Conditional~\cite{GeigerRH2026} \\
CRM/QG & Capacity + composition & $\tanh$ theorem~\cite{GeigerFST2026} \\
\textbf{Hodge} & \textbf{GHR spectrum} & \textbf{This paper} \\
\hline
\end{tabular}
\end{center}

In each case, the strategy is not to construct the desired object but to show that its non-existence would violate a positivity condition. The Hodge Conjecture is the last major instance where this strategy has not been applied systematically.

\subsection{Limitations}

\begin{enumerate}
\item \textbf{The hard direction is open.} This paper does not prove the Hodge Conjecture. It reformulates it as a positivity statement and identifies the precise obstruction, but the proof that arithmetic positivity forces algebraicity remains the central challenge.

\item \textbf{Partial sketches remain in Section~\ref{sec:functoriality}.} The K\"unneth functoriality is fully proved for primitive classes (Proposition~\ref{prop:kunneth}, Lemmas~\ref{lem:primitive_product}--\ref{lem:sign_product}), but the extension to non-primitive classes and the pushforward for fibred morphisms remain sketches requiring careful treatment of cross-terms in the primitive decomposition.

\item \textbf{Conjecture~\ref{conj:absolute_implies_ap} is unproven.} The relationship between absoluteness and universal Hodge-Riemann positivity is not established in general.

\item \textbf{Dependence on deep results.} The argument for abelian varieties (Theorem~\ref{thm:abelian}, Step~5) relies on Faltings' theorem (Tate conjecture for abelian varieties over finite fields) \cite{Faltings1983} and the Cattani--Deligne--Kaplan theorem on the algebraicity locus \cite{CattaniDeligneKaplan1995}.
\end{enumerate}

\subsection{Outlook}

Four directions are immediate:

\begin{enumerate}
\item \textbf{Verify the GHR Obstruction Conjecture for known non-algebraic classes.} The simplest test case is Voisin's counterexample to the integral Hodge conjecture \cite{Voisin2002b}: verify that her non-algebraic integral $(2,2)$-class on a product of K3 surfaces has negative GHR spectrum for some $\sigma$.

\item \textbf{Prove Conjecture~\ref{conj:absolute_implies_ap} for abelian varieties} using the explicit structure of the Mumford-Tate group and the symplectic Hodge-Riemann form.

\item \textbf{Develop the category of arithmetically positive classes} as a substitute for (part of) the theory of motives, with functorial operations and compatibility with the standard conjectures.

\item \textbf{The structure of the AP locus.} The set $\mathrm{AP}^p(X)$ is defined by \emph{quadratic} inequalities (the Hodge-Riemann form is bilinear, so the positivity conditions are quadratic in $\alpha$). It is therefore not a convex cone in the usual sense, but rather a \emph{quadratic cone} -- the intersection of a family of quadratic half-spaces indexed by polarisations and automorphisms. The set is closed under positive scaling ($\alpha \in \mathrm{AP}^p \Rightarrow \lambda\alpha \in \mathrm{AP}^p$ for $\lambda > 0$, since $Q_L(\lambda\alpha, \lambda\alpha) = \lambda^2 Q_L(\alpha, \alpha)$) and symmetric ($\alpha \in \mathrm{AP}^p \Rightarrow -\alpha \in \mathrm{AP}^p$). Understanding the boundary structure of this quadratic locus -- in particular, which classes lie on the boundary where $Q_L^\sigma(\alpha) = 0$ for some $\sigma$ and $L$ -- may provide geometric insight into the hard direction of the APC.
\end{enumerate}


% ===================================================================
% 10. CONCLUSION
% ===================================================================
\section{Conclusion}
\label{sec:conclusion}

We have proposed a reformulation of the Hodge Conjecture as a positivity No-Go theorem. The key concept -- \textit{arithmetic positivity} -- combines Hodge-Riemann positivity, absolute stability under $\mathrm{Aut}(\mathbb{C})$, and Lefschetz compatibility into a single condition that algebraic classes satisfy and (conjecturally) non-algebraic classes violate.

The approach offers three advantages over constructive methods:

\begin{enumerate}
\item \textbf{No cycle construction required.} The conjecture is reduced to a statement about positivity of bilinear forms, not the existence of subvarieties.

\item \textbf{No complete motivic machinery required.} The arithmetically positive cone is defined using only standard Hodge theory and Galois actions.

\item \textbf{The obstruction is explicit.} The Galois-Hodge-Riemann spectrum provides a computable (in principle) invariant that distinguishes algebraic from non-algebraic classes.
\end{enumerate}

The remaining challenge -- proving that non-negative GHR spectrum implies algebraicity -- is a precise, well-defined mathematical problem. We believe that the interplay between Hodge-Riemann positivity and Galois arithmetic provides the correct framework for attacking the Hodge Conjecture, and that the positivity No-Go perspective may succeed where constructive approaches have failed.

\begin{quote}
\textit{``The Hodge Conjecture does not ask us to build algebraic cycles. It asks us to show that geometry has no room for anything else.''}
\end{quote}


% ===================================================================
% ACKNOWLEDGMENTS
% ===================================================================
\begin{acknowledgments}
This work was conducted as independent research without institutional funding.

\textit{AI tools.} AI-based tools (Claude by Anthropic) were used for mathematical formalization and text generation. All mathematical content, conjectures, and responsibility for correctness lie solely with the author.

\textit{Funding information.} This research received no external funding.
\end{acknowledgments}


% ===================================================================
% BIBLIOGRAPHY
% ===================================================================
\begin{thebibliography}{30}

\bibitem{Voisin2002}
C.~Voisin (2002).
\textit{Hodge Theory and Complex Algebraic Geometry I, II}.
Cambridge University Press.
DOI: 10.1017/CBO9780511615344.

\bibitem{Deligne1982}
P.~Deligne (1982).
Hodge cycles on abelian varieties (Notes by J.~S.~Milne).
In \textit{Hodge Cycles, Motives, and Shimura Varieties}, Lecture Notes in Mathematics 900, pp.\ 9--100. Springer.
DOI: 10.1007/978-3-540-38955-2\_3.

\bibitem{Andre1996}
Y.~Andr\'e (1996).
Pour une th\'eorie inconditionnelle des motifs.
\textit{Publications Math\'ematiques de l'IH\'ES}, 83, 5--49.
DOI: 10.1007/BF02698643.

\bibitem{Charles2013}
F.~Charles (2013).
The Tate conjecture for K3 surfaces over finite fields.
\textit{Inventiones Mathematicae}, 194(1), 119--145.
DOI: 10.1007/s00222-012-0443-y.

\bibitem{Faltings1983}
G.~Faltings (1983).
Endlichkeitss\"atze f\"ur abelsche Variet\"aten \"uber Zahlk\"orpern.
\textit{Inventiones Mathematicae}, 73, 349--366.
DOI: 10.1007/BF01388432.

\bibitem{Abdulali1994}
S.~Abdulali (1994).
Algebraic cycles in families of abelian varieties.
\textit{Canadian Journal of Mathematics}, 46(6), 1121--1134.
DOI: 10.4153/CJM-1994-063-0.

\bibitem{Moonen1999}
B.~Moonen (1999).
Notes on Mumford-Tate groups.
Preprint, University of Amsterdam.

\bibitem{Voisin2002b}
C.~Voisin (2002).
A counterexample to the Hodge conjecture extended to K\"ahler varieties.
\textit{International Mathematics Research Notices}, 2002(20), 1057--1075.
DOI: 10.1155/S1073792802111135.

\bibitem{CattaniDeligneKaplan1995}
E.~Cattani, P.~Deligne \& S.~Kaplan (1995).
On the locus of Hodge classes.
\textit{Journal of the American Mathematical Society}, 8(2), 483--506.
DOI: 10.2307/2152824.

\bibitem{Grothendieck1969}
A.~Grothendieck (1969).
Standard conjectures on algebraic cycles.
In \textit{Algebraic Geometry (Bombay, 1968)}, pp.\ 193--199. Oxford University Press.

\bibitem{GrossZagier1986}
B.~Gross \& D.~Zagier (1986).
Heegner points and derivatives of $L$-series.
\textit{Inventiones Mathematicae}, 84, 225--320.
DOI: 10.1007/BF01388809.

\bibitem{Deligne1974}
P.~Deligne (1974).
La conjecture de Weil: I.
\textit{Publications Math\'ematiques de l'IH\'ES}, 43, 273--307.
DOI: 10.1007/BF02684373.

\bibitem{GeigerFST2026}
L.~Geiger (2026).
The Renormalized Free Energy Framework: A Unified Resolution of Structural Bottlenecks.
Zenodo preprint.
\href{https://doi.org/10.5281/zenodo.19036191}{doi:10.5281/zenodo.19036191}.

\bibitem{GeigerRH2026}
L.~Geiger (2026).
The Riemann Hypothesis I--III.
Zenodo preprint.
\href{https://doi.org/10.5281/zenodo.19035845}{doi:10.5281/zenodo.19035845}.

\end{thebibliography}

\end{document}
